% from ega0-19-fr.tex
\oldpage[0\textsubscript{IV-1}]{110}

\begin{theorem}[19.8.2]
\label{0.19.8.2-fr}
Soient $A$ un anneau local noethérien, $k$ son corps résiduel.
\begin{enumerate}
  \item[(i)] Si $B$ est une $A$-algèbre de Cohen, $B$ est une $A$-algèbre
  formellement lisse. Pour tout anneau local noethérien complet $C$, tout
  homomorphisme local $A\to C$ et tout idéal $\mathfrak J\neq C$ dans $C$,
  tout $A$-homomorphisme $B\to C/\mathfrak J$ se factorise donc en
  $B\xrightarrow{v}C\to C/\mathfrak J$, où $v$ est un $A$-homomorphisme
  (nécessairement local).
  \item[(ii)] Pour tout corps $K$, extension séparable de $k$, il existe une
  $A$-algèbre de Cohen $B$ telle que $B\otimes_A k$ soit $k$-isomorphe à $K$,
  et une telle $A$-algèbre est unique à isomorphisme près.
\end{enumerate}
\end{theorem}

\begin{proof}
Comme $K$ est une $k$-algèbre formellement lisse (19.6.1), l'assertion (i)
résulte de (19.7.1). Pour prouver (ii), on peut se borner au cas où $A$ est
\emph{complet}, car il revient au même de dire que $B$ est un $A$-module plat
ou un $\widehat A$-module plat ($\mathbf 0_{\mathrm{III}}$, 10.2.3), on a
$\mathfrak mB=\mathfrak m\widehat AB$ et $k$ est le corps résiduel de
$\widehat A$. Il suffit alors d'appliquer (19.7.2) en prenant
$\mathfrak J=\mathfrak m$ et $B_0=K$ (et utilisant (19.6.1)).
\end{proof}

\begin{definition}[19.8.3]
\label{0.19.8.3-fr}
On appelle \emph{anneau local premier} un anneau local de la forme
$\mathbf Z_{p\mathbf Z}$, où $p\mathbf Z$ est un idéal premier de $\mathbf Z$.
On appelle \emph{anneau local complet premier} le complété d'un anneau local
premier.
\end{definition}

Les anneaux locaux premiers sont donc de deux sortes~:
\begin{enumerate}
  \item[$1^{\mathrm o}$] Ceux qui correspondent aux idéaux maximaux
  $p\mathbf Z$ où $p\neq0$ est un nombre premier~; $\mathbf Z_{p\mathbf Z}$
  est un anneau de valuation discrète, dont le complété est \emph{l'anneau des
  entiers $p$-adiques} noté d'ordinaire $\mathbf Z_p$\footnote{Cette notation,
  actuellement universellement utilisée, est dans ce cas en conflit avec la
  notation $A_f$ adoptée dans ($\mathbf 0_{\mathrm I}$, 1.2.3)~: avec
  $A=\mathbf Z$ et $f=p$, $A_f$ signifie en effet l'anneau des nombres
  rationnels de la forme $k/p^n$ ($k\in\mathbf Z$, $n$ entier $\geq0$)~; nous
  éviterons toujours d'utiliser la notation $\mathbf Z_p$ pour désigner ce
  dernier anneau.}.
  \item[$2^{\mathrm o}$] Pour l'idéal premier $p\mathbf Z=(0)$,
  $\mathbf Z_{p\mathbf Z}$ est le corps des nombres rationnels $\mathbf Q$,
  identique à son complété (la topologie étant naturellement la topologie
  d'anneau local noethérien, donc ici la topologie discrète).
\end{enumerate}

La terminologie de (19.8.3), analogue à celle des «~corps premiers~», se
justifie de la même manière~: pour tout anneau local $A$, considérons
l'homomorphisme canonique $\varphi:\mathbf Z\to A$, et soit $p\mathbf Z$
l'image réciproque par cet homomorphisme de l'idéal maximal $\mathfrak m$ de
$A$~; $p\mathbf Z$ est un idéal premier de $\mathbf Z$ et l'homomorphisme
précédent se factorise donc en
$\mathbf Z\to\mathbf Z_{p\mathbf Z}\xrightarrow{\psi}A$~; d'ailleurs, comme
$\varphi$ est l'\emph{unique} homomorphisme de $\mathbf Z$ dans $A$, $p$ et
$\psi$ sont déterminés de façon \emph{unique}. Autrement dit, pour tout anneau
local $A$, il y a un \emph{unique homomorphisme} $\psi:P\to A$, où $P$ est un
anneau local premier~; si en outre $A$ est séparé et complet, on peut prolonger
par complétion cet homomorphisme, et il y a donc un \emph{unique
homomorphisme} $\widehat\psi:\widehat P\to A$, où $\widehat P$ est un anneau
local complet premier. D'ailleurs, par passage aux quotients, $\psi$ donne un
homomorphisme du corps résiduel $\mathbf F_p$ si $p>0$ (resp. $\mathbf Q$ si
$p=0$) dans le corps résiduel $k$ de $A$, et $p$ est donc la
\emph{caractéristique} de $k$.

Si l'on prend en particulier pour $A$ un anneau local premier (resp. complet
premier), on voit qu'il n'existe dans un tel anneau qu'\emph{un seul
endomorphisme}, savoir l'identité.
